Napjainkban rendkívül elterjedté váltak a webes alkalmazások, amelyeket nap mint nap használunk: akár a munkánk során, akár kapcsolattartásra, akár szórakozásra. Ez a hatalmas növekedés a webes alkalmazások használatában annak köszönhető, hogy egy böngészőből használható alkalmazás már-már az asztali programok színvonalát és funkcionalitását is meg tudja közelíteni.
A modern webes alkalmazások mögött rengeteg tervezés, döntés, fejlesztés, valamint technikai vagy akár pénzügyi akadály áll, amelyek megoldásáért és kivitelezéséért egy csapat felelős. Mindenkinek megvan a szerepe az ötlet kigondolásától a UI (felhasználói felület) tervezésén át egészen a fejlesztésig és karbantartásig.
A dolgozat egy ilyen webes alkalmazás tervezési folyamatait kívánja bemutatni: hogy miként jutunk el magától az ötlettől egy kész termékig, és milyen folyamatok mennek végbe még azelőtt, hogy a fejlesztők leírnának egy sor kódot, vagy a UX-designerek akár egy vonalat is húznának.
Ennek demonstrálásához bemutatok egy „fiktív” fejlesztői csapatot, akik egy filmnaplózó weboldalt terveznek és fejlesztenek.

\textbf{Kulcsszavak}: webes alkalmazások, tervezés, fejlesztés, UX-design, filmnaplózó